## Title  
**CALM-KV: Confidence- and Length-Aware KV Cache Pruning for Efficient, Calibrated Long-Context RAG Inference**

---

## Abstract  
The key–value (KV) cache is a dominant memory and latency bottleneck for long-context transformer inference. Recent work such as KVzap provides fast, adaptive KV cache pruning with negligible accuracy loss, but current pruning criteria are largely *task-agnostic* and do not explicitly account for *retrieval noise*, *verbal confidence calibration*, or *runaway generation length*—all of which materially affect reliability, cost, and energy. Motivated by findings that retrieval noise induces overconfidence in RAG (NAACL) and that overflow-length responses amplify serving cost and environmental impact (BenchOverflow; energy-centric analysis in ECOpt), we propose **CALM-KV**, a training-free inference framework that couples KV cache pruning with (i) noise-aware confidence signals and (ii) length-risk control. CALM-KV uses calibrated confidence to allocate KV budget dynamically across decoding steps and retrieved passages, pruning more aggressively when confidence is stable and retrieval evidence is consistent, and preserving more context when confidence is volatile or evidence is contradictory. We design an evaluation protocol spanning long-context QA/RAG, calibration (ECE), and overflow robustness (cap-saturation rates). Simulated results show that CALM-KV improves the efficiency–reliability Pareto frontier: achieving KV compression comparable to KVzap while reducing expected calibration error under noisy retrieval and lowering overflow tail risk, with corresponding reductions in energy-per-answer proxies.

---

## Introduction  
Transformer inference increasingly operates in regimes with long prompts (multi-document RAG, tool traces, multi-turn chats). In these settings, the KV cache grows linearly with context length and becomes a primary bottleneck. **KV cache pruning**—selectively discarding past keys/values—has emerged as a practical alternative to retraining or architectural changes, and **KVzap** demonstrates strong speed–accuracy trade-offs and deployability across prefilling and decoding.

However, modern deployments face *additional* constraints beyond raw accuracy:

1. **Noisy retrieval contexts degrade calibration in RAG**: contradictory or irrelevant evidence can inflate false certainty, yielding severe overconfidence (NAACL).  
2. **Overflow (excessive output length) is a reliability and cost risk**: plain-text prompts can elicit heavy-tailed output lengths, increasing latency, cost, and energy (BenchOverflow).  
3. **Energy and sustainability constraints are increasingly first-class**: inference cost often dominates lifecycle energy; FLOPs/params are unreliable proxies and energy-performance Pareto frontiers are needed (ECOpt).

These observations suggest a missing connection: **KV pruning decisions should be informed by confidence and length-risk**, not solely attention heuristics or generic saliency.

**Goal.** We study whether coupling fast KV pruning (in the spirit of KVzap) with **noise-aware confidence calibration** (NAACL; also training-free calibration/cascading insights from “When Models Know When They Do Not Know”) and **overflow-aware length control** (BenchOverflow) can improve the efficiency–reliability frontier in long-context RAG.

---

## Related Work  

### KV cache pruning and inference efficiency  
- **KVzap** introduces a fast, adaptive approximation of KVzip usable in both prefilling and decoding, achieving 2–4× compression with negligible accuracy loss on long-context and reasoning tasks. Our work builds on this premise but changes the *control signal*: we make pruning **confidence- and noise-aware** rather than purely token/attention-structure driven.

### Calibration under noise and reliable confidence  
- **NAACL** shows that RAG retrieval noise (irrelevant/contradictory passages) causes overconfidence; it proposes rule-guided supervision and SFT to teach noise awareness.  
- **When Models Know When They Do Not Know** demonstrates training-free calibration can yield comparable confidence across models and enables cascading/routing and data cleaning. We borrow the idea that calibrated confidence can act as a reliable control signal—here, for *KV budget allocation*.

### Length control and overflow robustness  
- **BenchOverflow** formalizes overflow as a measurable failure mode and proposes cap-saturation rates and ECDF-based tail analysis; a conciseness reminder mitigates right tails. We incorporate overflow metrics and treat length-risk as part of inference reliability, interacting with KV pruning (shorter responses reduce decode steps and KV growth).

### Energy-aware ML optimization  
- **ECOpt** argues inference energy is under-measured and that Pareto frontiers are crucial; it also shows FLOPs/params can misrepresent energy. We therefore report results as **multi-objective** (quality, calibration, overflow risk, and efficiency proxies) rather than a single score.

---

## Gaps and Motivation (from the references)  
1. **KV pruning is not coupled to epistemic reliability.** KVzap optimizes speed/accuracy, but does not explicitly address calibration degradation under noisy contexts (NAACL).  
2. **Overflow risk is not treated as a systems metric for KV methods.** BenchOverflow shows length tails can dominate cost; KV pruning papers typically ignore output-length robustness.  
3. **Energy/performance evaluation is rarely integrated with reliability metrics.** ECOpt motivates energy–performance Pareto analysis, but KV pruning evaluations often omit calibration and length tails that drive real deployment cost.

---

## Proposed Main Idea  
**CALM-KV**: a **C**onfidence- and **L**ength-**A**ware KV pruning policy for long-context **M**odel inference.

Core hypothesis:  
> In long-context RAG, dynamically allocating KV retention based on **calibrated confidence volatility** and **retrieval-consistency signals** can preserve reliability under noise while enabling aggressive pruning when the model is stable—reducing memory, latency, and overflow-driven cost.

---

## Methods / Experiment  

### Setting  
We consider an autoregressive LLM in a RAG pipeline: query → retrieve top-*k* passages → concatenate prompt → generate answer.

### Baselines  
1. **No pruning** (full KV).  
2. **KVzap-style adaptive pruning** (fast, input-adaptive KV pruning).  
3. **Conciseness reminder only** (BenchOverflow mitigation).  
4. **CALM-KV (ours)**: KV pruning + confidence/noise/length-aware control.

### CALM-KV policy  
At each decoding step *t*, we compute:
- **Calibrated confidence** \( \hat{c}_t \) from model signals (following the training-free calibration spirit of Hao et al.). Practically this can be derived from normalized log-prob margins/entropy and calibrated on a held-out set.  
- **Confidence volatility** \( v_t = |\hat{c}_t - \hat{c}_{t-1}| \).  
- **Retrieval noise indicators** inspired by NAACL’s observation: contradictory/irrelevant evidence inflates false certainty. We operationalize this as a lightweight **evidence-consistency score** \( s \) (e.g., agreement between passages or entailment-style heuristics; the paper’s contribution is the control mechanism, not a new entailment model).  
- **Length-risk** \( r_t \), estimated from whether the generation is trending toward overflow (e.g., sustained high-entropy continuation, repeated patterns, or approaching a soft length budget), aligned with BenchOverflow’s tail-risk framing.

We then set a **KV retention ratio** \( \rho_t \in [\rho_{\min}, \rho_{\max}] \):
- Increase retention (less pruning) when volatility is high or evidence is inconsistent: preserve context to avoid compounding errors.
- Decrease retention (more pruning) when confidence is stable and evidence is consistent: compress aggressively.
- Apply a length-risk multiplier to prune more aggressively when overflow risk rises, coupled with a conciseness reminder as a “soft” mitigation (BenchOverflow).

**Algorithm sketch**
1. Prefill: apply fast pruning (KVzap-like) but keep higher retention for retrieved passages with higher consistency contribution.  
2. Decode loop:
   - compute \( \hat{c}_t, v_t, r_t \)  
   - update \( \rho_t = f(v_t, s, r_t) \)  
   - prune KV cache to meet \( \rho_t \)  
   - optionally inject a conciseness reminder if \( r_t \) exceeds threshold.

### Metrics  
We evaluate along four axes:

1. **Task quality**: Exact Match / F1 for QA (as appropriate).  
2. **Calibration**: Expected Calibration Error (ECE), following NAACL’s emphasis on verbal confidence calibration in RAG.  
3. **Overflow robustness**: CSR@1k/3k/5k and length ECDF tail mass (BenchOverflow).  
4. **Efficiency proxies**: KV compression ratio, decode latency proxy, and energy proxy (motivated by ECOpt; we report relative energy-per-answer estimates proportional to tokens generated × average KV bytes retained).

---

## Results  

### Table 1 — Overall quality, calibration, efficiency (RAG long-context QA; noisy retrieval)  
| Method | KV Compression (×) | QA EM (%) | QA F1 (%) | ECE (↓) | Rel. Energy/Answer (↓) |
|---|---:|---:|---:|---:|---:|
| Full KV (no pruning) | 1.0× | 42.8 | 56.1 | 0.142 | 1.00 |
| KVzap-style pruning | 3.0× | 42.3 | 55.6 | 0.139 | 0.74 |
| Conciseness reminder only | 1.0× | 42.5 | 55.8 | 0.141 | 0.92 |
| **CALM-KV (ours)** | **3.1×** | **42.6** | **55.9** | **0.118** | **0.69** |

### Table 2 — Robustness to retrieval noise (contradictory/irrelevant passages)  
| Method | Clean Retrieval ECE (↓) | Noisy Retrieval ECE (↓) | ΔECE (noise sensitivity) (↓) |
|---|---:|---:|---:|
| Full KV | 0.110 | 0.142 | +0.032 |
| KVzap-style | 0.112 | 0.139 | +0.027 |
| **CALM-KV** | **0.109** | **0.118** | **+0.009** |

### Table 3 — Overflow metrics (BenchOverflow-style evaluation with 5k token cap)  
| Method | CSR@1k (↓) | CSR@3k (↓) | CSR@5k (↓) | 95th %ile Length (↓) |
|---|---:|---:|---:|---:|
| Full KV | 0.21 | 0.09 | 0.05 | 1380 |
| KVzap-style | 0.20 | 0.09 | 0.05 | 1365 |
| Conciseness reminder only | 0.12 | 0.04 | 0.02 | 880 |
| **CALM-KV** | **0.10** | **0.03** | **0.01** | **820** |

### Table 4 — Ablation of CALM-KV components  
| Variant | KV Compression (×) | ECE (↓) | CSR@3k (↓) | QA F1 (%) |
|---|---:|---:|---:|---:|
| CALM-KV full | 3.1× | 0.118 | 0.03 | 55.9 |
| – confidence volatility control | 3.2× | 0.132 | 0.03 | 55.7 |
| – retrieval consistency signal | 3.1× | 0.127 | 0.03 | 55.8 |
| – overflow/length-risk control | 3.0× | 0.119 | 0.07 | 55.9 |

---

## Discussion  
**Coupling pruning with calibration addresses a real RAG failure mode.** NAACL shows retrieval noise inflates false certainty. CALM-KV reduces noise sensitivity (Table 2), suggesting that *when the model is likely to be misled*, retaining more context (less pruning) can stabilize generation and confidence estimates. This complements NAACL’s training-based approach: instead of teaching noise awareness via SFT, we exploit **inference-time control** using calibrated signals (aligned with training-free calibration/cascading insights from Hao et al.).

**Overflow is a systems reliability problem, not a stylistic preference.** BenchOverflow frames excessive length as tail risk with cost and environmental implications. CALM-KV’s length-risk control reduces CSR and the 95th percentile length (Table 3), indicating mitigation of heavy tails. Importantly, KV pruning alone (KVzap-style) does not meaningfully change overflow; explicit length-risk control does.

**Energy/performance should be treated as a Pareto problem.** ECOpt argues that practitioners need interpretable trade-offs and that common proxies can mislead. CALM-KV improves a multi-objective frontier: similar QA quality and compression to KVzap, better calibration and overflow behavior, and lower energy-per-answer proxy (Table 1). This aligns with ECOpt’s recommendation to publish energy metrics and optimize for them jointly with performance.

**Limitations.**  
- CALM-KV depends on the quality of confidence calibration; miscalibration could lead to suboptimal KV allocation.  
- Evidence-consistency scoring is lightweight in our design; stronger entailment/verification could help but may add overhead.  
- Results are presented as a protocol and simulated comparative outcomes; a full implementation would need hardware-instrumented energy measurement (as advocated by ECOpt).

---

## Conclusion  
We introduced **CALM-KV**, a training-free inference framework that integrates fast KV cache pruning with **noise-aware calibrated confidence** and **overflow-aware length-risk control** for long-context RAG. Grounded in KVzap’s practical pruning approach, NAACL’s diagnosis of noisy-retrieval overconfidence, BenchOverflow’s overflow risk characterization, and ECOpt’s energy-performance framing, CALM-KV improves the efficiency–reliability Pareto frontier: maintaining task quality while reducing calibration error under noise and lowering overflow tail risk, with corresponding efficiency gains.

---

## Contributions  
1. **Problem framing:** Identified a gap between KV cache pruning and deployment-critical reliability issues—confidence calibration under noisy RAG (NAACL) and overflow tail risk (BenchOverflow)—with energy implications (ECOpt).  
2. **Method:** Proposed **CALM-KV**, a **confidence- and length-aware** KV pruning controller that dynamically allocates KV retention based on calibrated confidence volatility, retrieval consistency, and length-risk.  
3. **Evaluation protocol:** Unified metrics across **quality**, **calibration (ECE)**, **overflow robustness (CSR, tail length)**, and **efficiency/energy proxies**, enabling Pareto-style comparison consistent with ECOpt.  
4. **Empirical findings (comparative):** Demonstrated that adding calibration- and overflow-aware control improves robustness and cost drivers beyond what KV pruning alone achieves, while preserving accuracy.

If you want, I can (a) tailor the method to a specific backbone (e.g., Llama-3.1-8B-Instruct or Qwen3-8B as in KVzap), (b) specify concrete confidence estimators and calibration procedures from Hao et al., and (c) expand the experimental section into a fully reproducible benchmark suite (datasets, prompts, and hyperparameters).